\section{Other caps and angular partitions}
\label{ext:cap-section}

The triangular construction leaves two choices available: how to end a
sector radially, and where to sample its direction.  The same balance
between radial enlargement and perpendicular projection permits curved
caps, smaller fixed margins, and sectors centred at the sampling nodes.
We continue to use
\[
 H=\sqrt2,\qquad
 K=\{(x,y,Hs):0\le s\le1,\ |x|\le1-s,\ |y|\le s\},
 \qquad d(q)=\frac{q^2+q^4}{16}.
\]
Thus $A_{\min}(K)=H$.  All the cylinders below are specified by an
intercept set and a nonzero direction; their bases are the perpendicular
projections of the intercept sets.

\subsection{A fixed positive margin}

The following criterion isolates what the angular construction uses.
Continuity of the displaced boundaries chooses a sector, and a positive
quadratic margin ensures that the two chosen radial caps cannot both fail.
The partition may be nonuniform, and its number of intervals need not
satisfy an upper bound.

\begin{theorem}[Compatible boundaries and radial caps]
\label{ext:general-cap}
Let $\tau>0$ tend to zero through any set accumulating at zero.
For each $\tau$, choose a finite partition
$-1=q_0<\cdots<q_m=1$, constants $\alpha_j,\beta_j$, and positive
continuous functions $R_j$ on $I_j=[q_{j-1},q_j]$.
Suppose that $L(q)=\alpha_j-q\beta_j$ is continuous on $[-1,1]$
and $L(-1)=L(1)=0$.  Suppose also, uniformly for every $j$ and
$q\in I_j$, that
\begin{align}
 \beta_j&=q/2+O(\tau^2),&
 \alpha_j&=(1+q^2)/4+O(\tau^2),\label{ext:coefficient-hypotheses}\\
 R_j(q)&=\tfrac12+\tau^2\bigl(d(q)+\eta\bigr)+O(\tau^4),
 &&0<\eta<1/480.\label{ext:cap-hypothesis}
\end{align}
Here $\eta$ is fixed, and all error constants and their common range of
validity are independent of $\tau,j,q$.
Define
\begin{align*}
 P_j^-&=\{(0,qS,HS):q\in I_j,\ 0\le S\le R_j(q)\},&
 v_j^-&=(1,\tau\alpha_j,H\tau\beta_j),\\
 P_j^+&=\{(qT,0,H(1-T)):q\in I_j,\ 0\le T\le R_j(q)\},&
 v_j^+&=(-\tau\alpha_j,1,H\tau\beta_j).
\end{align*}
For every sufficiently small admissible $\tau$, the $2m$ cylinders
$P_j^-+\mathbb Rv_j^-$ and $P_j^++\mathbb Rv_j^+$ cover the closed
tetrahedron $K$.  Their bases are compact, and their total area $C_\tau$
satisfies
\begin{equation}\label{ext:general-cost}
 \frac{C_\tau}{H}
 =\frac12+\left(2\eta-\frac1{240}\right)\tau^2+O(\tau^4).
\end{equation}
If each $R_j$ is constant, every base is a triangle.
The threshold and error constant may depend on $\eta$ and the uniform
constants in the hypotheses, but not on $m$.
\end{theorem}

\begin{proof}
Fix $(x,y,Hs)\in K$.  The continuous piecewise affine function
$q\mapsto qs+\tau xL(q)$ has outer values $-s,s$, which bracket $y$.
Choose the first partition node after $q_0$ whose value is at least $y$.
Its preceding value is at most $y$, by the initial endpoint inequality
or by minimality.  On the selected interval $I_j$, put
\[
 S=s-\tau x\beta_j,\qquad Y=y-\tau x\alpha_j.
\]
The crossing inequalities say $q_{j-1}S\le Y\le q_jS$, so $S\ge0$.
If $S=0$, then $Y=0$ and the selected cylinder already covers the point.
Otherwise $Y=qS$ for some $q\in I_j$.
Applying the same argument to
$p\mapsto p(1-s)-\tau yL(p)$ gives an interval $I_i$ and
\[
 T=1-s+\tau y\beta_i\ge0,\qquad
 X=x+\tau y\alpha_i=pT,\qquad p\in I_i.
\]
Again $T=0$ already gives coverage.  It remains to exclude
$S>R_j(q)$ and $T>R_i(p)$ simultaneously.

Assume these failures.  For sufficiently small $\tau$, the fixed
positive margin and the uniform cap error imply $S,T>1/2$.
Since $S+T=1+\tau(y\beta_i-x\beta_j)=1+O(\tau)$, each equals
$1/2+O(\tau)$.  The intercept equations then give
$y=q/2+O(\tau)$ and $x=p/2+O(\tau)$.
Write $e_j=\beta_j-q/2$ and $e_i=\beta_i-p/2$.
The same complementary radial balance as in the triangular construction
now has explicit sampling errors:
\begin{align*}
 s-\tau xy
 &=\tfrac12+(S-\tfrac12)(1-\tau xq)
       -\tau^2x^2\alpha_j+\tau xe_j,\\
 1-s+\tau xy
 &=\tfrac12+(T-\tfrac12)(1+\tau yp)
       -\tau^2y^2\alpha_i-\tau ye_i.
\end{align*}
These identities follow by substituting
$y=qS+\tau x\alpha_j$ and $x=pT-\tau y\alpha_i$.
Their left sides sum to one, and both factors multiplying the radial
excesses are positive for small $\tau$.  Substitute the strict cap
failures, divide by $\tau^2$, and use $e_i,e_j=O(\tau^2)$ to obtain
\[
 0> d(q)+d(p)+2\eta-x^2\alpha_j-y^2\alpha_i+O(\tau)
   =2\eta+\frac{(p^2-q^2)^2}{16}+O(\tau).
\]
The error is uniform.  This is impossible for sufficiently small
$\tau$, proving coverage, including the boundary.

The parametrizations of $P_j^-$ and $P_j^+$ have area Jacobians $HS$
and $HT$, respectively, and are one-to-one off their tips.
Projection multiplies area in either intercept plane by
$[1+\tau^2(\alpha_j^2+2\beta_j^2)]^{-1/2}$.
Consequently
\[
 \frac{C_\tau}{H}=
 \sum_j\int_{I_j}\frac{R_j(q)^2}
 {\sqrt{1+\tau^2(\alpha_j^2+2\beta_j^2)}}\,dq.
\]
Uniform expansion of the integrand gives
\[
 \frac14+\tau^2\left[
 d(q)+\eta-\frac18\left(\frac{(1+q^2)^2}{16}+\frac{q^2}{2}\right)
 \right]+O(\tau^4).
\]
The total angular length is two, so integration preserves the uniform
fourth-order remainder regardless of the number of intervals.  Finally
\[
 \int_{-1}^1d(q)\,dq=\frac1{15},\qquad
 \frac18\int_{-1}^1\left(\frac{(1+q^2)^2}{16}+\frac{q^2}{2}\right)dq
 =\frac{17}{240},
\]
which proves \eqref{ext:general-cost}.  Each parameter domain is compact;
its continuous image and perpendicular projection are compact as well.
Projection is invertible between the intercept plane and the base plane,
so a constant cap gives a nondegenerate triangle.
\end{proof}

\subsection{Triangular and curved apertures}
\label{ext:apertures}

For an ordinary interval $I=[l,r]$, use
\begin{equation}\label{ext:ordinary-secants}
 \alpha_I=\frac{1+lr}{4},\qquad \beta_I=\frac{l+r}{4}.
\end{equation}
The identity $\alpha_I-q\beta_I=(1-q^2)/4$ at both endpoints
matches adjacent boundaries exactly.  If the maximum interval length
is at most $M\tau^2$, for fixed $M$, then
\eqref{ext:coefficient-hypotheses} holds uniformly.
With $r_I=(l+r)/2$, each of the following choices satisfies
\eqref{ext:cap-hypothesis}:
\begin{align}
 R_I(q)&=\tfrac12+\tau^2\bigl(d(r_I)+\eta\bigr),
 &&\text{midpoint triangles},\label{ext:midpoint-cap}\\
 R_I(q)&=\tfrac12+\tau^2\bigl(\max_{u\in I}d(u)+\eta\bigr),
 &&\text{maximum triangles},\label{ext:maximum-cap}\\
 R_I(q)&=\tfrac12+\tau^2\bigl(d(q)+\eta\bigr),
 &&\text{continuous radial caps},\label{ext:radial-cap}\\
 R_I(q)&=\sqrt{\tfrac14+\tau^2\bigl(d(q)+\eta\bigr)},
 &&\text{squared-radius caps}.\label{ext:squared-cap}
\end{align}
Indeed, $d$ has bounded derivative on $[-1,1]$, so replacing $d(q)$
by its midpoint value or its cell maximum costs $O(\tau^2)$ inside
the parentheses.  For \eqref{ext:squared-cap}, Taylor expansion at
$1/4$ gives the required fourth-order error.  Thus all four choices
have the area formula \eqref{ext:general-cost}.  The first two give
triangles; the last two give compact radial sectors.

Another choice specifies the radial enlargement through the transverse
intercept.  Put $D(Y)=Y^2/4+Y^4$.  For $b=3/4$ or $b=1$, retain the
bounded aperture
\begin{equation}\label{ext:cartesian-aperture}
 0\le S\le b,\qquad S\le\tfrac12+\tau^2\bigl(D(qS)+\eta\bigr).
\end{equation}
For small $\tau$, this is exactly $0\le S\le R(q)$, where $R$ is
continuous and satisfies \eqref{ext:cap-hypothesis}.  To verify this,
consider $f(S)=S-1/2-\tau^2(D(qS)+\eta)$ on $[0,1]$.  Uniformly in
$|q|\le1$,
\[
 f'(S)=1-\tau^2(q^2S/2+4q^4S^3)\ge1-\tfrac92\tau^2.
\]
For small $\tau$, $f$ is strictly increasing, is negative at $1/2$,
and is positive at $3/4$.  Its unique zero $R(q)$ lies between these
values and depends continuously on $q$.  Its defining equation first
gives $R(q)-1/2=O(\tau^2)$, and then
\[
 D(qR(q))=D(q/2)+O(\tau^2)=d(q)+O(\tau^2),
\]
which proves the cap expansion.  The outer bound in
\eqref{ext:cartesian-aperture} is part of the definition: without it,
the quartic inequality can admit a second, unbounded component.

The common conclusion is easiest to compare in the fixed tetrahedron $K$:
for each fixed $0<\eta<1/480$, every cap above has normalized total area
\eqref{ext:general-cost}.  Its threshold is uniform over partitions whose
maximum interval length is at most $M\tau^2$, for a fixed $M$.
For example, an equal partition with any integer $m\ge2/\tau^2$ gives
exactly $2m$ cylinders, with the same threshold and error constant as the
mesh is refined.  Each cover is finite; no further limit in the number
of cylinders is needed after choosing the tilt and mesh.

The squared-radius and Cartesian formulas also explain the geometry in
two different ways.  Appendix~\ref{ext:squared-direct} compares the two
radial squares, while Appendix~\ref{ext:cartesian-direct} uses a
discriminant to estimate the intercept height without locating the
perturbed interface.  These are alternative proofs of cases already
covered by Theorem~\ref{ext:general-cap}.  Exact counts and costs for
several rescaled parameter choices are collected in
Appendix~\ref{ext:scaling-specializations}.

\subsection{Sampling at the endpoints}
\label{ext:endpoint-sectors}

The compatibility condition also permits directions sampled at nodes
rather than along secants of each sector.  Let $k>1$ be an integer,
$\tau=1/k$, $n=k^2$, and $q_j=-1+2j/n$ for $0\le j\le n$.
Put $\lambda_0=-1$, $\lambda_{n+1}=1$, and
$\lambda_j=(q_{j-1}+q_j)/2$ for $1\le j\le n$.
Use the sectors $[\lambda_j,\lambda_{j+1}]$ and parameters
\[
 \beta_j=q_j/2,\qquad \alpha_j=(1+q_j^2)/4,\qquad
 R_j=\tfrac12+\tau^2\bigl(d(q_j)+1/1000\bigr).
\]
At an interior boundary,
$\alpha_j-\alpha_{j-1}=\lambda_j(\beta_j-\beta_{j-1})$;
at the exterior boundaries $\alpha_j-\lambda\beta_j=0$.
Thus the piecewise affine boundary function is continuous and vanishes
at both ends.  Every angular point of the $j$th sector is within $\tau^2$
of $q_j$, so all the hypotheses of Theorem~\ref{ext:general-cap} hold.
There are $n+1$ sectors, with widths $\tau^2$ at the two ends and
$2\tau^2$ in the interior.  Consequently, for all sufficiently large $k$,
this construction covers $K$ with exactly $2(k^2+1)$ triangular bases
and normalized total area
\begin{equation}\label{ext:endpoint-cost}
 \frac{C_\tau}{A_{\min}(K)}
 =\frac12-\frac{13}{6000}\tau^2+O(\tau^4).
\end{equation}
The half-width end sectors preserve the exact exterior boundaries.
Appendix~\ref{ext:endpoint-direct} gives a complementary radial identity
that verifies their coverage directly, and
Appendix~\ref{ext:endpoint-scaled} records the corresponding absolute
area at edge length $\sqrt2$.
